% Discussion and interpretation of results
% for direct inclusion in a LaTeX paper

\section{Discussion}\label{sec:discussion}

\subsection{Evidence Type Asymmetry at G5}

The central finding of this analysis is the evidence type asymmetry at
the V\&V goal (G5). Four evidence types converge at this node
(Table~\ref{tab:evidence-types}), each differing in three dimensions:
what is measured (code structure, scenario space, model confidence,
attacker capability), how the evidence is produced (static analysis,
simulation, inference, red-teaming), and what scale the result uses
(binary pass/fail, scenario count, statistical distribution, attack
success rate).

This asymmetry creates a concrete challenge for the practitioner: how
to determine that V\&V evidence is \emph{sufficient} when four
incommensurable evidence types must jointly support the claim. ISO~26262
prescribes MC/DC coverage at ASIL~D as a binary pass criterion.
ISO~21448 prescribes triggering condition coverage as a count with a
project-specific threshold. ISO/PAS~8800 prescribes uncertainty
evaluation but defines no numerical threshold. ISO/SAE~21434 prescribes
penetration testing but leaves the acceptable attack success rate to the
organisation.

No standard defines a combination rule, a weighting scheme, or a
minimum set of evidence types that together constitute sufficiency.
This gap is structural: it arises from the independent development of
each standard by different technical committees (ISO/TC~22/SC~32 for
ISO~26262, ISO/TC~22/SC~32/WG~14 for ISO~21448, ISO/TC~22/SC~32/WG~11
and SAE for ISO/SAE~21434, ISO/TC~22/SC~32/WG~18 for ISO/PAS~8800).

\subsection{Integration-Induced Gaps}

Two of the six identified gaps (Gap-3 and Gap-4) are
\emph{integration-induced}: they are invisible when each standard is
applied independently and only become apparent when the integrated
argument is constructed.

Gap-3 (adversarial-SOTIF boundary) arises because ISO~21448 explicitly
excludes cybersecurity threats (Clause~1), while ISO/SAE~21434 includes
adversarial scenarios. An adversarial point cloud perturbation that causes
a false negative simultaneously satisfies the definition of a SOTIF
triggering condition (a condition that leads to a hazardous event
without a system fault) and a cybersecurity threat scenario (an
intentional attack exploiting a vulnerability). The boundary between
G7 and G8 in the integrated GSN is undefined, and no standard provides
guidance for this ownership decision.

Gap-4 (no cross-domain release decision criteria) arises because the
integration produces four separate evidence sets, each assessed against
its own acceptance framework. ISO~26262-2 Clause~6.4 prescribes a
functional safety assessment; ISO~21448 Clause~12 prescribes a SOTIF
release decision; ISO/SAE~21434 Clause~3.1.11 prescribes a cybersecurity
case evaluation. No standard prescribes how to make a single release
decision that considers all four evidence sets jointly.

\subsection{Quantitative Asymmetry}

The integrated argument reveals a quantitative asymmetry between
hardware-oriented and AI-oriented assurance claims. ISO~26262 Part~5
prescribes specific quantitative targets for hardware: single-point
fault metric (SPFM~$\geq$~99\% at ASIL~D), latent fault metric
(LFM~$\geq$~90\% at ASIL~D), and probabilistic metric for random
hardware failures (PMHF~$<$~$10^{-8}$~h$^{-1}$). No equivalent
quantitative targets exist for AI component reliability.

The closest analogue is the SOTIF acceptance criterion (ISO~21448
Clause~6.5), which is qualitative and project-specific. ISO/PAS~8800
does not define reliability targets. ISO/IEC~TR~5469 acknowledges the
problem (Clause~9.2.2: non-separability of system behaviour) but does
not provide targets. This creates Gap-1 in the classification.

\subsection{Evaluation Results}

The CARLA evaluation demonstrates the operational consequence of the
SOTIF triggering conditions identified in the specification phase.
Under non-triggering conditions, mean recall remains above the baseline
performance level. Under triggering conditions (rain~$>$~20~mm/h or
visibility~$<$~200~m), mean recall drops and ensemble geometric
divergence increases, indicating elevated model uncertainty.

The divergence metric serves dual roles in the integrated GSN:
(i)~as V\&V evidence at G5 (ISO/PAS~8800 Clauses~8--9), demonstrating
that the deep ensemble quantifies its own uncertainty under adverse
conditions, and (ii)~as a runtime monitoring indicator at G6
(ISO/PAS~8800 Clause~14), where elevated divergence signals potential
out-of-distribution inputs. This dual role illustrates how a single
technical mechanism can provide evidence for multiple nodes in the
integrated argument.

\subsection{Implications for Practitioners}

The identified inconsistencies and gaps have practical consequences for
automotive engineers constructing assurance cases for AI-based perception
components:

\begin{enumerate}
    \item \textbf{Evidence combination}: Practitioners must define a
          project-specific combination rule for the four evidence types
          at G5, as no standard provides one. This rule should be
          documented as an explicit assumption in the assurance case.
    \item \textbf{Boundary definition}: The adversarial-SOTIF boundary
          (Gap-3) must be resolved at the organisation level. A
          recommended approach is to assign ownership based on the
          attacker model: scenarios with a defined adversary belong to
          ISO/SAE~21434; scenarios caused by environmental conditions
          without an adversary belong to ISO~21448; scenarios where both
          apply require coordinated assessment.
    \item \textbf{Integrated release}: The release decision (Gap-4) should
          be structured as a single gate that evaluates all four evidence
          sets against defined criteria, rather than four independent gates.
    \item \textbf{Modification re-assurance}: Until G9 is developed by
          future standards work, OTA updates to AI models should trigger
          a re-evaluation of all affected claims in the integrated
          argument, not just the claims from the standard most directly
          related to the change.
\end{enumerate}

\subsection{Limitations}

This analysis has three primary limitations.
First, the integrated GSN is constructed from normative requirements and
does not account for organisation-specific interpretations or tailoring
permitted by each standard.
Second, the CARLA evaluation uses synthetic results generated from a
parametric degradation model, not from a trained SECOND detector deployed
in a live CARLA simulation. The evaluation demonstrates the methodology
and evidence structure, not the performance of a specific model.
Third, the analysis covers four normative standards and one informative
TR. Other potentially relevant standards (UL~4600, UNECE R155/R156,
EU AI Act) are not included.
